\section{The forced-hit gate calculus}
\label{sec:gates}

The ordinary zone charges every future defender placement before a role or
window deadline.  A tight forcing line contains replies that the defender
must copy into the certificate or abandon the line.  An exact copied reply
does not create a real-only frontier stone.  This permits a narrower hazard
clock.  The escape branch still needs the full remaining defender turn.
The construction and coordinate checks in this section are recorded in the
companion proof and forced-hit report \cite{hexobot-companions}.

\subsection{Why pressure alone cannot be debited}

The first example isolates the terminal race.  The current threats appear to
consume both defender placements, but the defender can ignore them and win
before the attacker receives a turn.

\begin{example}[Reachable disjoint-hit race]
Let
\[
\begin{aligned}
W&=\{(0,r):0\le r\le5\},\\
T_1&=\{(q,-4):0\le q\le5\}, &a&=(5,-4),\\
T_2&=\{(q,8):-5\le q\le0\}, &b&=(-5,8).
\end{aligned}
\]
The following legal history begins with a defender opening at the origin.
Each row after the opening is one two-placement turn.
\begin{center}
\small
\begin{tabular}{@{}ccl@{}}
\toprule
Plies & Mover & Placements \\
\midrule
$0$ & $\D$ & $(0,0)$ \\
$1$--$2$ & $\A$ & $(0,-4),(1,-4)$ \\
$3$--$4$ & $\D$ & $(0,1),(-1,-4)$ \\
$5$--$6$ & $\A$ & $(2,-4),(3,-4)$ \\
$7$--$8$ & $\D$ & $(0,2),(0,-8)$ \\
$9$--$10$ & $\A$ & $(4,-4),(-4,8)$ \\
$11$--$12$ & $\D$ & $(0,3),(1,8)$ \\
$13$--$14$ & $\A$ & $(-3,8),(-2,8)$ \\
$15$--$16$ & $\D$ & $(-8,0),(8,-8)$ \\
$17$--$18$ & $\A$ & $(-1,8),(0,8)$ \\
\bottomrule
\end{tabular}
\end{center}
Every placement is within distance eight of an earlier stone.  No displayed
prefix completes a window.  The resulting node has defender budget $b=2$ and
complete attacker threat-empty family
\[
\mathcal F=\bigl\{\{a\},\{a,(6,-4)\},
                    \{b\},\{b,(-6,8)\}\bigr\},
\qquad \tau(\mathcal F)=2.
\]
The singleton sets force every transversal to contain $a$ and $b$, so the
extendable-hit kernel is exactly $\{a,b\}$.  It is disjoint from $W$.
The defender may ignore both threats and play
\[
u=(0,4),\qquad v=(0,5).
\]
Both cells are legal at turn start.  The first raises $\cnt_\D(W)$ to five,
and the second completes $W$.  Termination occurs before the attacker can use
either surviving threat.
\end{example}

The same construction can be translated or reflected, and remote legal
fillers can be added.  It is not a boundary artifact.  A sound debit must
distinguish a copied hit from a refusal and must charge the refusal branch for
the complete remaining turn.  Threat masks also have to agree between the
real and ghost positions at the gate.  Otherwise a real-only defender stone
can kill one named threat and reduce the actual pressure.

Substitution raises a separate obstruction.  A real hit and a different
ghost hit can answer the same threats while creating nonempty $X$ and $Y$
sets.  The real stone may then support a radius-eight relay that the ghost
stone does not support.  Such a hit remains an ordinary transition for the
frontier clock.  Only an exact copied placement receives the gate debit.

\subsection{Protected exact-copy gates}

A gate names the current pressure and protects its masks through an entry
checkpoint.  It then searches an exact residual-transversal kernel rather
than an ordinary four-part zone.

\begin{definition}[Protected exact-copy forcing gate]
\label{def:D19}
At an internal AND node $Q$ with defender budget $b\in\{1,2\}$, a
\emph{forcing gate} names a finite family $\mathcal H_Q$ of current attacker
threat windows.  Put
\[
\mathcal F_Q=\{E(U,P_Q):U\in\mathcal H_Q\}.
\]
The verifier checks
\[
\tau(\mathcal F_Q)=b,
\qquad \neg\ownwinnow(P_Q).
\tag{FG1--FG2}
\]
For a cell $d$, define
\[
\mathcal F_Q\setminus d=\{E\in\mathcal F_Q:d\notin E\},
\qquad
K(Q)=\{d\in\Legal(P_Q):\tau(\mathcal F_Q\setminus d)\le b-1\}.
\]
For every $U\in\mathcal H_Q$ and $e\in E(U,P_Q)$, add a checkpoint role
$(Q,U,e)$ to the obligation grammar.  Its deadline is the gate-entry mask
check immediately before the defender reply.  It belongs to every strict
ancestor's reachable obligation union and remains protected through that
check.  Write $\Prot^-(Q)$ for the incoming protected set.  After the
verifier checks the named ghost masks, the checkpoint roles are discharged
and the remaining set is $\Prot^+(Q)$.

The searched set at the gate is exactly the finite nonempty set $K(Q)$.  Each
$d\in K(Q)$ has its exact successor $C_d$, and a continuing real reply is
copied on that edge.  No substitution envelope is permitted on a kernel
reply.  A real reply outside $K(Q)$ abandons the original subtree and uses
the adaptive escape contract of \cref{lem:L15}.

Let $p(Q)$ be the absolute index of the last placement before entry to $Q$.
The escape deadline is
\[
p(Q)+b+2.
\]
The certificate horizon is the maximum of its ordinary declared resolutions
and all reachable escape deadlines.
\end{definition}

\begin{plainterms}
The checkpoint makes every named threat the same in the real and ghost
positions before the gate is used.  A reply inside the kernel is copied at
the same cell.  A reply outside it discards the old plan and gives the
attacker a surviving current threat, unless the defender wins first during
the remainder of the turn.
\end{plainterms}

\begin{figure}[tbp]
  \centering
  \begin{tikzpicture}[node distance=14mm and 14mm]
    \node[cert gate] (q) {$Q$: gate\\$\tau(\mathcal F_Q)=b$,
      $\neg\ownwinnow$\\checkpoint masks equal};
    \node[cert node,below left=of q] (c1) {$C_{d_1}$\\exact child};
    \node[cert node,below=of q] (c2) {$C_{d_2}$\\exact child};
    \node[cert node,below right=of q] (esc) {old subtree abandoned\\
      adaptive surviving threat};
    \draw[cert edge] (q) -- node[left,font=\scriptsize] {$d_1\in K(Q)$} (c1);
    \draw[cert edge] (q) -- node[right,font=\scriptsize] {$d_2\in K(Q)$} (c2);
    \draw[cert escape] (q) -- node[right,font=\scriptsize,align=left]
      {$d\notin K(Q)$\\deadline $p(Q)+b+2$} (esc);
    \node[font=\scriptsize,align=center,below=4mm of esc]
      {attacker completes unless defender\\wins during the current turn};
  \end{tikzpicture}
  \caption{Forcing-gate grammar.  Kernel edges are exact copies.  The dashed
  escape is an adaptive terminal contract, not a substituted certificate
  child.}
  \label{fig:forcing-gate}
\end{figure}

The kernel is nonempty.  If $T$ is a size-$b$ minimum transversal, each
$d\in T$ leaves $T\setminus\{d\}$ as a transversal of
$\mathcal F_Q\setminus d$, and \cref{thm:T1} makes $d$ legal.  It is finite
because a kernel cell must belong to some named empty set.  At $b=1$, one
threat is enough.  At $b=2$, a subfamily of at most three threats can retain
transversal number two.

The role clock counts opportunities to create a real-only defender stone.
An exact gate edge creates none.

\begin{definition}[Forced-hit-debited role ranks]
\label{def:D20}
For a role $\rho$ live at $N$, define $f_N(\rho)$ in reverse topological
order by
\[
f_N(\rho)=
\begin{cases}
0,&\rho\text{ is at its deadline or no longer reachable},\\
f_C(\rho),&N\text{ is an ordinary OR and }\rho\text{ remains live},\\
1+\max_C f_C(\rho),&N\text{ is an ordinary AND node},\\
\max_{d\in K(N)}f_{C_d}(\rho),&N\text{ is a forcing gate}.
\end{cases}
\]
A child where $\rho$ is not reachable contributes zero.  For a cell carrying
several live roles, set
\[
f_N(y)=\max\{f_N(\rho):\rho\text{ is live at }N\text{ and carried by }y\}.
\]
An ordinary AND edge costs one.  An exact gate edge costs zero because its
defender stone is shared.  An off-kernel reply abandons every old role.
\end{definition}

\begin{plainterms}
The value $f$ is not the number of defender plies before resolution.  It
counts only placements that can extend a real-only legality chain toward a
protected cell.  Full path length continues to use the local budget $B$.
\end{plainterms}

A copied gate hit can still harm a defender window when the hit lies inside
that window.  The exposure recurrence therefore keeps a window-specific
indicator and a full escape floor.

\begin{definition}[Forced-hit-debited window exposure]
\label{def:D20a}
For a forcing gate, write
\[
G_N(W)=\max_{d\in K(N)}
  \bigl(\mathbb{1}[d\in W]+Q_{C_d}^\D(W)\bigr).
\]
For every window $W$, define $Q_N^\D(W)$ by
\[
Q_N^\D(W)=
\begin{cases}
0,&N\text{ is WIN or OR-COMPLETION},\\
b,&N\text{ is a LOSS leaf with defender budget }b,\\
0,&N\text{ is an OR whose attacker placement enters }W,\\
Q_C^\D(W),&N\text{ is any other ordinary OR},\\
1+\max_C Q_C^\D(W),&N\text{ is an ordinary AND},\\
\max\{b,G_N(W)\},&N\text{ is a forcing gate}.
\end{cases}
\]
If $W$ is already not alive for the defender, set $Q_N^\D(W)=0$.  The
maximum at a gate is taken child by child.  Replacing it by independent
maxima is conservative but can overcount by one.

On a continuing path, let $F$ count every ordinary defender opportunity,
every remaining defender placement at a LOSS leaf, and all $b$ placements of
the first off-kernel escape turn.  Let $H_W$ count exact copied gate replies
whose cells lie in $W$.  The recurrence is the branch-coherent maximum of
$F+H_W$.  An ordinary reply is counted in $F$ even when it hits a threat.
\end{definition}

\begin{plainterms}
A copied hit outside $W$ costs zero to that window.  A copied hit inside $W$
costs one.  Refusing a gate costs the full remaining turn, and a LOSS leaf
also costs its full remaining budget.  These cases prevent the race from the
opening example.
\end{plainterms}

The ordinary four-part zone now uses the debited hazard clocks.  A gate itself
does not use these four terms.

\begin{definition}[Ordinary debited zone]
\label{def:D21}
At an ordinary internal AND node, define
\[
\begin{aligned}
\ZdirFH(N)&=\Prot(N)\cap\Legal(P_N),\\
\ZseedFH(N)&=
 \bigcup_{\substack{y\in\Prot(N)\setminus
 (\Legal(P_N)\cup\Stones(P_N))\\ f_N(y)\ge1}}
 \left(\Legal(P_N)\cap\Ball{8(f_N(y)-1)}{\{y\}}\right),\\
\ZtouchFH(N)&=\bigcup\left\{E(W,P_N):
 \begin{array}{l}W\text{ is defender-alive},\ \cnt_\D(W,P_N)\ge1,\\
 \cnt_\D(W,P_N)+Q_N^\D(W)\ge6
 \end{array}\right\},\\
\ZvirginFH(N)&=\left\{c\in\Legal(P_N):
 \begin{array}{l}\text{some all-empty }W\text{ has }Q_N^\D(W)\ge6,\\
 \dist(c,W)\le8(Q_N^\D(W)-6)
 \end{array}\right\}.
\end{aligned}
\]
The node searches an independently nonempty superset of the union of these
four sets.  Its obligation union includes every reachable checkpoint role,
and clause \textup{(Z4)} is unchanged.  A forcing gate performs its entry
checks and searches exactly $K(Q)$.

For a fixed augmented certificate,
\[
f_N(\rho)\le r_N(\rho),
\qquad Q_N^\D(W)\le E_N^\D(W).
\]
The full $B$, $r$, and $E^\D$ clocks retain their ordinary AND charges at a
gate.  Checkpoint roles can enlarge $\Prot$ at strict ancestors, so no claim
is made that the total searched-set size decreases in every certificate.
\end{definition}

\begin{plainterms}
The debit can shrink a ranked seed band or a completion guard for a fixed
augmented certificate.  It does not make the gate an ordinary zone node, and
it does not discount a general substitution.  Added checkpoint obligations
can offset the smaller hazard radii.
\end{plainterms}

\subsection{Transfer at a gate}

The entry checkpoint makes the named threat masks identical.  The residual
transversal inequality then covers every reply outside the kernel.

\begin{lemma}[Protected gate dichotomy]
\label{lem:L15}
On entry to a \cref{def:D19} gate, every named threat window has identical
real and ghost masks.  For every real legal defender reply $d$, exactly one
of the following holds.
\begin{enumerate}[label=\textup{(\arabic*)}]
  \item If $d\in K(Q)$, then $d$ is shared-empty and ghost-legal.  Both games
  place at $d$, take the exact child, and leave $X$ and $Y$ unchanged.
  \item If $d\notin K(Q)$ and the defender does not win during the remaining
  current turn, the attacker completes a named surviving threat in at most
  two placements of the following turn, by $p(Q)+b+2$.
\end{enumerate}
\end{lemma}

\begin{proof}
Fix $U\in\mathcal H_Q$.  Its attacker stones are shared.  Every ghost empty
$e\in E(U,P_Q)$ carries a checkpoint role, so $e\notin X$ at the entry
check.  Since $U$ is attacker-alive in the ghost position, it contains no
ghost defender stone and hence no $Y$-cell.  The complete real and ghost
masks of $U$ agree.

If $d\in K(Q)$, then $d$ belongs to a named empty set.  It is shared-empty and
is legal by \cref{thm:T1}.  The gate supplies its exact child.  This is the
first case.

Suppose $d\notin K(Q)$.  By definition,
\[
\tau(\mathcal F_Q\setminus d)>b-1.
\]
Let $H$ be the set of at most $b-1$ later defender placements in the current
turn.  It cannot meet every member of $\mathcal F_Q\setminus d$.  Some named
window therefore avoids both $d$ and $H$.  It remains attacker-alive and
retains its one or two initial empties.  By \cref{lem:L1}, those empties lie
within distance two of permanent shared attacker stones.  They are legal, and
the attacker completes the window in the next turn.  The last possible
attacker placement has index $p(Q)+b+2$.  The game rule allows the defender
to terminate first during the current turn, which is the alternative retained
in the statement.  Membership in $K(Q)$ makes the two cases disjoint.
\end{proof}

The weighted clocks bound the two hazards needed by the coupling.  They also
remain below the corresponding full clocks.

\begin{lemma}[Weighted hazard bounds]
\label{lem:L16}
For every certificate using \cref{def:D19,def:D20,def:D20a,def:D21}:
\begin{enumerate}[label=\textup{(\arabic*)}]
  \item on a continuing mapped path, the number of ordinary real-only
  frontier opportunities before a live role's deadline is at most
  $f_N(\rho)$;
  \item for a fixed $W$, before the certificate attacker wins or enters $W$,
  or before an off-kernel escape resolves, the weighted count of ordinary
  defender edges, copied gate edges in $W$, and all placements in the first
  LOSS or escape remainder is at most $Q_N^\D(W)$;
  \item $f_N(\rho)\le r_N(\rho)$ and $Q_N^\D(W)\le E_N^\D(W)$;
  \item at a gate, $B(Q)\ge b$, so the full local budget covers every escape
  remainder.
\end{enumerate}
\end{lemma}

\begin{proof}[Proof sketch]
Reverse induction follows the displayed recurrences.  An ordinary AND edge
can add one $X$-stone and is charged once; a copied gate edge changes neither
$X$ nor $Y$ and is charged to $W$ exactly when its cell lies in $W$.  LOSS
and escape remainders retain the full base $b$.  Comparing each clause with
the ordinary rank and exposure recurrences gives the two inequalities, while
the nonempty kernel and the $b=2$ successor phase give $B(Q)\ge b$.
See the companion proof for the child-by-child induction
\cite{hexobot-companions}.
\end{proof}

The first-bad-event argument combines the frontier and completion channels.

\begin{lemma}[Debited first-bad-event lemma]
\label{lem:L17}
Under \cref{def:D19,def:D20,def:D20a,def:D21}:
\begin{enumerate}[label=\textup{(\arabic*)}]
  \item while the real play remains mapped to the certificate, no defender
  placement creates a real-only stone in the current protected set; old roles
  are abandoned when an off-kernel escape starts;
  \item before the mapped certificate attacker or a gate-escape attacker
  resolves, no real defender play completes a window.
\end{enumerate}
\end{lemma}

\begin{proof}[Proof sketch]
For protected occupation, trace a first bad cell back to the last ghost-legal
dismissed seed.  Copied gate stones are shared and cannot be internal links of
a real-only chain, so every link is charged by $f$ and the seed lay in
$\ZseedFH$.  For a first defender completion, separate no-dismissal, touched,
and virgin cases.  The mask identity handles no-dismissal lines; $Q^\D$
charges every ordinary fill, every copied hit in the fixed window, and each
full LOSS or escape remainder.  The touched count or virgin distance
inequality then places the first dismissed fill or its seed in the applicable
zone.  The complete first-event analysis is in the companion proof
\cite{hexobot-companions}.
\end{proof}

The gate theorem replaces the two ordinary hazard clocks while retaining the
rest of the certificate grammar.

\begin{theorem}[Exact-copy $F+H_W$ soundness]
\label{thm:T11}
Let a finite certificate tree satisfying \cref{def:D9}, or a finite
certificate DAG satisfying \cref{def:D18}, use \cref{def:D21} at every
ordinary AND node and \cref{def:D19} at every forcing gate.  Include every
checkpoint role in the reachable obligation unions, retain clause
\textup{(Z4)}, retain the full local budget $B$ of \cref{def:D14}, and include
every gate escape deadline in the global horizon.  Then the compiled attacker
wins against every real defender play by that horizon.  The debited $f$ and
$Q^\D$ values are sound replacements for $r$ and $E^\D$ in the ordinary
debited seed, touched, and virgin terms.
\end{theorem}

\begin{proof}
At an ordinary node, run the A1--A3 coupling from \cref{thm:T3}.  Use
\cref{def:D21,lem:L17} in place of the ordinary protected-occupation and
completion lemmas.

At a gate, check the incoming masks and discharge the checkpoint roles.
Lemma~\ref{lem:L15} transfers the named threat family.  A real reply in
$K(Q)$ is shared-empty, is copied on its exact child, and leaves $X$ and $Y$
unchanged.  A reply outside $K(Q)$ abandons the mapped subtree.  The second
part of \cref{lem:L17} excludes a defender completion during the remaining
turn.  The second case of \cref{lem:L15} then supplies an adaptive surviving
threat that the attacker completes by the escape deadline.

WIN, LOSS, OR-COMPLETION, and ordinary OR transfers are those of
\cref{thm:T3}.  A LOSS remainder remains sound because $Q^\D$ charges all
$b$ placements and the obligation grammar protects every witness empty
through leaf entry.  A finite tree reaches a typed terminal or its first
escape.  A finite DAG can be unfolded as in \cref{thm:T10}; the gate labels,
checkpoint roles, and branch-coherent maxima are preserved.  Hence every real
play resolves by the augmented horizon.
\end{proof}

\begin{plainterms}
The theorem permits a discount only after the gate masks have been protected
and the real reply is copied at the same cell.  Every refusal is charged as an
escape, and every ordinary or substituted reply keeps its full transition
cost.  These conditions preserve the main coupling while making the two
hazard clocks target-specific.
\end{plainterms}

An ordinary substitution envelope can coexist with gates, but its transition
continues to use full clocks.

\begin{theorem}[Substitution-envelope compatibility]
\label{thm:T11.1}
The conclusion of \cref{thm:T11} remains valid when an ordinary node's global
debited-zone tests are replaced by a valid \cref{def:D17} envelope, provided
the selected reachable-role union includes every future checkpoint role and
the envelope retains the transition-inclusive $\widehat B$, role ranks, and
$\widehat E$ tests of \cref{def:D17}.  No debited substitution envelope is
asserted.
\end{theorem}

\begin{plainterms}
An ordinary substitution still creates a real--ghost mismatch on its current
ply, so it receives no gate discount.  Its full transition clocks must also
protect every later gate checkpoint.  Omitting either condition would combine
two individually sound reductions without covering their interaction.
\end{plainterms}

\begin{proof}[Proof sketch]
The real reply and ghost substitute create the ordinary $X,Y$ transition, so
the current unit remains mandatory.  The full envelope covers all three
causal channels and every later checkpoint role.  Subsequent debited ordinary
steps satisfy $f\le r$ and $Q^\D\le E^\D$, while a gate retains the full
envelope inequalities and charges any escape floor.  The nested-envelope
argument of \cref{thm:T9} then applies \cite{hexobot-companions}.
\end{proof}

\subsection{The running certificate fires its gate}

The position introduced in \cref{sec:certificates} has defender budget one,
a singleton threat with empty cell $h=(5,0)$, and kernel $K=\{h\}$.  The
gate copies $h$.  The attacker then plays $(3,20)$ and $(0,20)$ and reaches a
defender-budget-two LOSS leaf whose three witness-empty pairs are disjoint.
Their transversal number is three.

For the defender-alive window
\[
W=\{(q,40):0\le q\le5\},
\]
the root count is three and $h\notin W$.  The full exposure is
\[
E_{\mathrm{root}}^\D(W)=1+2=3,
\]
so its touched test reads $3+3=6$.  The gate recurrence instead gives
\[
Q_{\mathrm{root}}^\D(W)=\max\{1,0+2\}=2,
\]
and the debited test reads $3+2=5$.  The value two is attained when the LOSS
remainder places at $(3,40)$ and $(4,40)$.  The window reaches count five,
one of the three LOSS witnesses survives, and the attacker completes that
witness.  The one-unit debit and the LOSS base two are both used on a legal
line.

For
\[
W'=\{(q,0):5\le q\le10\},
\]
the copied hit lies in the queried window.  The later LOSS placements
$(6,0)$ and $(7,0)$ give
\[
Q_{\mathrm{root}}^\D(W')=\max\{1,1+2\}=3.
\]
All three placements lie in $W'$ and the count does not reach six.  Omitting
the indicator for a copied hit inside $W'$ would give two and would undercount
this line.

\subsection{What remains open}

The exact-copy theorem is proved, but it resolves only the two named hazards
at protected tight gates.  The broader debit objective remains partial.

\begin{table}[tbp]
\centering
\small
\begin{tabular}{@{}p{.24\linewidth}p{.68\linewidth}@{}}
\toprule
Open extension & Present boundary \\
\midrule
Target-independent budget & No scalar $B_{FH}<B$ is proved.  Path length,
LOSS remainders, substitution envelopes, and escape remainders still use the
full $B$. \\
Substituted forced hits & No D17 substitute receives zero cost.  A different
real and ghost hit can create a real-only frontier corridor. \\
Net searched-set shrinkage & The $f$ and $Q^\D$ terms do not grow for a fixed
augmented certificate, but checkpoint roles can enlarge ancestor obligation
sets. \\
Slack pressure & A family with $\tau<b$ does not force the next reply; its
kernel is the whole legal set.  No additional debit is proved for that turn. \\
Seed-radius sharpness & Exact sharpness of every $f$-based reduction remains
open.  Ordinary paths retain the relative boundary $8(r-1)$. \\
\bottomrule
\end{tabular}
\caption{Open extensions of the forced-hit gate calculus.}
\label{tab:gate-gaps}
\end{table}

The limit map in \cref{sec:tightness} records this result as
\emph{partially resolved}.  That status concerns the intended pruning
extension, not the validity of \cref{thm:T11}.
